% ============================================================
% PROJECTS — full (14 entries)
% Source: section_projets.tex (deprecated)
% Original used \project macro inside \begin{multicols}{2}\begin{projects}.
% Converted to tabularx two-column style.
% ============================================================

\subsection{SELECTED RESEARCH PROJECTS}
\vspace{2pt}

\noindent
\begin{tabularx}{\textwidth}{@{}lX@{}}
2026--2033 & \textbf{ARC Centre of Excellence for Quality Work in a Digital Age (Qwida)} \\
           & \emph{Chief Investigator, Proposal co-author} \\
           & The Centre is composed of 7 Australian universities. The Centre's vision is to enable a future where workers and organisations are empowered by technology, rather than beholden to it. \\
           & \textit{University of Melbourne, 94M AUD} \\[3pt]

2026--2030 & \textbf{Balancing Autonomy and Control in Interactive Social Agent Design} \\
           & \emph{Project Lead, Proposal co-author} \\
           & This ARC Future Fellowship develops frameworks for designing AI-driven social agents that balance autonomy with user control. It introduces adaptive control and feedback systems for socially aligned, safe, and inclusive human-AI interactions across education, aged care, and remote collaboration contexts. \\
           & \textit{Lead PI, Australian Research Council, FT250100459, 4 Years Fellowship} \\[3pt]

2026--2029 & \textbf{Robotic Trainers for a Skilled Workforce of the Future} \\
           & \emph{CI, Proposal co-author} \\
           & This ARC Discovery Project develops robotic tools that act as adaptive trainers to accelerate human skill acquisition. It models task difficulty, designs multimodal feedback, and integrates AI-based adaptation to guide and assess trainees, addressing Australia's skilled labour shortages. \\
           & \textit{Co-PI, Australian Research Council, DP260101082, 4 Years Project} \\[3pt]

2024--2025 & \textbf{Implication of Industry 4.0 Technologies on Work Practices} \\
           & \emph{Proposal co-author} \\
           & The project aims to provide insights and guidelines to inform the choice, design and integration of technology that maximizes desirable work practices and improve productivity and working conditions in the building and construction sectors. \\
           & \textit{Building 4.0 CRC, 433k AUD} \\[3pt]

2024--2025 & \textbf{A3: An Automatic Academic Advisor Tool using Data-Driven Curriculum Exploration for Designing Learning Pathways} \\
           & \emph{Project Lead, Proposal co-author} \\
           & The project aims to design and develop a solution for visualisation the curriculum by creating a knowledge map using data parsed from the LMS. \\
           & \textit{LTI 2023, University of Melbourne, 30k AUD} \\[3pt]

2024       & \textbf{Adaptive Robot Autonomy based on Human's Mental State: The Case of Cognitive Load and Trust} \\
           & \emph{Project Lead, Proposal co-author} \\
           & This project aims to develop user-centered tools that enhance the interaction between workers and robotic systems by adjusting the level of autonomy of the robot in response to the human's mental state (i.e.\ cognitive load but also human's understanding of the robot's behaviour). \\
           & \textit{CIS \& ME 2023 Research Collaboration Seed Funding, University of Melbourne, 30k AUD} \\[3pt]

2021--2024 & \textbf{Human-Centered Robot Learning} \\
           & \emph{Proposal author, Lead CI} \\
           & This project is funded by the ARC. The goal of the project is to investigate novel interaction methods to allow novices to teach robots basic tasks for at home deployment. \\
           & \textit{440k AUD, 3 Years Project, ARC DECRA} \\[3pt]

2021       & \textbf{GIRL -- Gamified Interactive Robot Learning} \\
           & \emph{Proposal author, Lead CI} \\
           & The project aims to investigate how gamification could help novices to engage in robot teaching activities. \\
           & \textit{UNSW Faculty of Engineering, 20k AUD, GROW Award} \\[3pt]

2019--2022 & \textbf{IReCheck} \\
           & \emph{Proposal co-author, International Partner} \\
           & The IReCheck project, is funded by the Swiss NSF and the French ANR. The goal of the project is to investigate social and task adaptation in a Learning by Teaching task with Typically Developed children and children with Neuro-Developmental Disorders. \\
           & \textit{CHILI EPFL, 490k CHF, 3 Years Project} \\[3pt]

2018--2022 & \textbf{Animatas} \\
           & \emph{Proposal co-author, International Partner} \quad \href{http://animatas.eu}{animatas.eu} \\
           & The Animatas project, is a EU MSCA ITN ETN that aims to explore the use of virtual agents in educational context. \\
           & \textit{CHILI EPFL, 759k CHF, 3 Years Project, 3 PhD Students} \\[3pt]

2018--2020 & \textbf{Cowriting Kazakh} \\
           & \emph{Principal Coordinator, Proposal co-author} \quad \href{https://chili.epfl.ch/cowriter}{chili.epfl.ch/cowriter} \\
           & The Cowriter Kazakh project, is a SNSF project that seeds for collaboration between Kazakhs and Swiss researchers. In this project we will deploy the Cowriter robotic system to help the transition from Cyrillic to Roman alphabet in Kazakhstan. \\
           & \textit{CHILI EPFL, 10k CHF, 2 Years Project, Nazarbayev University} \\[3pt]

2018--2020 & \textbf{Find your voice!} \\
           & \emph{Local Coordinator, Proposal co-author} \\
           & The project aims to help children who struggle to find their voice in social situations, by training them on joke telling using platform such as Alexa or Google Home. It is a project funded by the SRCD that seeds for cross-country and cross discipline collaboration between researchers. \\
           & \textit{CHILI EPFL, 15k CHF, Partners: UCLIC, University of Sussex, UC Irvine} \\[3pt]

2015--2019 & \textbf{Cellulo} \\
           & \emph{Project manager, Proposal co-author} \quad \href{http://chili.epfl.ch/cellulo}{chili.epfl.ch/cellulo} \\
           & The Cellulo project, funded by the NCCR Robotics (Swiss NSF), aims to design and propose new use of tangible swarm robots in education. \\
           & \textit{CHILI/MOBOTS EPFL, 1M CHF, 4 Years, 2 PhD + 1 Postdoc} \\[3pt]

2015--     & \textbf{Robots for Learning} \\
           & \emph{Project manager, Proposal co-author} \quad \href{https://robots4learning.wordpress.com/}{robots4learning.wordpress.com} \\
           & I launched in Fall 2015 the Robots for Learning initiative that aims to create a research community on interactions aspect of using robots for learners. We organized several workshop and training events funded by the NCCR Robotics. \\
           & \textit{CHILI/MOBOTS EPFL, 18.5k CHF} \\
\end{tabularx}
